\begin{abstract}
PROPHET predicts which of 18 attacker actions comes next from the
Windows audit records received before a prediction cutoff. Under benign
workloads absent from every labeled training run, models sharing data
and architecture diverge across seeds, and the final training-batch
loss did not rank the seeds. Consensus distillation (CD), a standard
ensemble-distillation setup, trains one student toward the averaged
probabilities of five baseline models on unlabeled inputs from the new
workload. In a sealed in-house evaluation on new runs, mean negative
log-likelihood (NLL) fell from 0.0731 to 0.0276. On the public ADS-24
dataset, in a development comparison on inspected runs, the seed
standard deviation fell from 0.163 to 0.025, yet every distilled seed
of our model stayed worse than a training-prior baseline (NLL 1.830).
Three general encoders moved the same way in development comparisons,
though not every seed improved. CD narrowed the spread among these
checkpoints without removing the loss under environment shift.
% 근거: RECORD:166, 170 (v3, A+C FLAT); RECORD:218--219, 221, 234 (ADS-24);
% RECORD:294 (세 인코더). 세 인코더 결과는 자체·ADS-24 모두 이미 열람한 실행의
% 개발 비교다(answers 2026-11-24 07:10, 4번). "20 to 30 of 30"은 seed별 원자료를
% atlas에서 찾지 못해 뺐다. "not every seed improved"는 9장의 검증된 사례(ADS-24
% LSTM 기준 seed 2개가 기준선 아래, 학생은 0개)가 받친다. "training-prior baseline" = FIT 실행 동등 전술 빈도 예측기(8장 정의).
\end{abstract}

\begin{Engkeyword}
Next attacker action prediction, Underspecification, Seed instability,
Consensus distillation, Unlabeled adaptation, Windows audit logs
\end{Engkeyword}

\begin{korabstract}
본 연구는 예측 시점 전에 수신된 윈도우 감사 로그만으로 공격자가 다음에
수행할 행동을 18개 범주 중에서 예측한다. 정답이 붙은 학습 실행에 없던 정상
작업 조건에서는 같은 자료와 구조로 학습한 모델도 시드마다 성능이 크게
갈리며, 학습 마지막 미니배치의 손실로는 시드의 순위를 가릴 수 없었다. 이
흔들림을 줄이기 위해 기존의 앙상블 증류 방식을 적용하여, 새 작업 조건의
정답 없는 입력에서 기준 모델 다섯 개가 합의한 확률을 학생 모델이 따라
배우도록 한다. 학생은 기준 모델과 구조가 같아 추론 때도 단일 모델이다.
사전에 봉인한 자체 실험실 평가의 새 실행에서 평균 음의 로그우도는
0.0731에서 0.0276으로 낮아졌다. 공개 자료 ADS-24에서는 이미 열람한 평가
실행에서의 개발 비교로 시드 간 표준편차가 0.163에서 0.025로 줄었다. 세 가지
일반 인코더도 개발 비교에서 같은 방향을 보였다. 다만 모든 시드가 좋아지지는
않았고, ADS-24에서 우리 모델의 증류 시드는 모두 학습 자료의 전술 빈도만 쓰는
기준선(1.830)보다 손실이 높았다. 이 방법은 학습한 시드들 사이의 흔들림을 줄이지만, 환경 이동에 따른
성능 저하 자체를 없애지는 못한다.
\end{korabstract}

\begin{Korkeyword}
공격 행동 예측, 과소명세, 시드 불안정성, 합의 증류, 무라벨 적응, 윈도우 감사 로그
\end{Korkeyword}
